\section{Произведение мер и теорема Фубини}
\begin{definition}
    Пусть $X$ — абстрактное множество. Семейство $\mathcal{M} \subset 2^X$ называется \textit{монотонным классом}, если:
    \begin{enumerate}
        \item $\{A_i\} \subset \mathcal{M}$ и $A_1 \subset A_2 \subset \dots \subset A_n \subset \dots$, то $\bigcup\limits_{i=1}^{\infty} A_i \in \mathcal{M}$,
        \item $\{B_j\} \subset \mathcal{M}$ и $B_1 \supset B_2 \supset \dots \supset B_n \supset \dots$, то $\bigcap\limits_{n=1}^{\infty}B_n \in \mathcal{M}$.
    \end{enumerate}
\end{definition} 

\begin{example}
    Любая $\sigma$-алгебра $\mathfrak{M} \subset 2^X$ является монотонным классом.
\end{example}

\begin{theorem}[О монотонном классе]
    Пусть $X$ — абстрактное множество.
    Пусть $\mathcal{A}$ — алгебра подмножеств множества $X$.
    Если $\mathcal{M} \supset \mathcal{A}$ — монотонный класс, то $\mathcal{M} \supset \mathfrak{M}(\mathcal{A})$ — минимальная $\sigma$-алгебра, порождённая $\mathcal{A}$.
\end{theorem} 
\begin{proof}
    Пусть $\mathcal{M}(\mathcal{A})$ — минимальный монотонный класс, содержащий алгебру $\mathcal{A}$.
    Заметим, что $\mathcal{M}(\mathcal{A}) \subset \mathfrak{M}(\mathcal{A})$, поскольку $\mathfrak{M}(\mathcal{A})$ сама является монотонным классом.
    Покажем, что на самом деле $\mathcal{M}(\mathcal{A}) = \mathfrak{M}(\mathcal{A})$.
    Рассмотрим при $A \in \mathcal{A}$ множество
    \[\mathcal{M}_A := \{B\in \mathcal{M}(\mathcal{A}) \ | \ A \cap B \in \mathcal{M}(\mathcal{A}) \text{ и } A \cap B^c \in \mathcal{M}(\mathcal{A})\}.\]
    Заметим, что поскольку $\mathcal{A} \subset \mathcal{M}(\mathcal{A})$, то $\mathcal{M}_A \supset \mathcal{A}$.
    С другой стороны, $\mathcal{M}_A$ — монотонный класс:
    \[A \cap \left(\bigcup\limits_{i=1}^{\infty}B_i\right)^c = \bigcap\limits_{i=1}^{\infty} A \cap B_i^c, \ B_1 \subset \dots \subset B_n \subset \dots,\]
    следовательно, $\mathcal{M}_A = \mathcal{M}(\mathcal{A})$ для любого $A \in \mathcal{A}$.
    
    В частности, можно взять $A = X$ и получить $\mathcal{M}_X = \mathcal{M}(\mathcal{A})$, следовательно, $\mathcal{M}(\mathcal{A})$ замкнут относительно взятия дополнения.
    Кроме того, мы получили, что для любого $A \in \mathcal{A}$ и для любого $B \in \mathcal{M}(\mathcal{A})$ верно $A \cap B \in \mathcal{M}(\mathcal{A})$ и $A \cap B^c \in \mathcal{M}(\mathcal{A})$.
    
    Теперь для любого $B \in \mathcal{M}(\mathcal{A})$ рассмотрим
    \[N_B := \{A \in \mathcal{M}(\mathcal{A}) \ | \ B \cap A \in \mathcal{M}(\mathcal{A})\}.\]
    Очевидно, что $\mathcal{A} \subset N_B$.
    Кроме того, $N_B$ является монотонным классом, откуда следует, что $N_B = \mathcal{M}(\mathcal{A})$.

    Мы доказали, что $\mathcal{M}(\mathcal{A})$ является симметричной (содержит множество и его дополнение) системой и замкнуто относительно конечных пересечений, также он по определению монотонного класса замнкут относительно счётного возрастающего объединения, следовательно, $\mathcal{M}(\mathcal{A})$ — $\sigma$-алгебра, а потому $\mathcal{M}(\mathcal{A}) = \mathfrak{M}(\mathcal{A})$.
\end{proof}

\paragraph{Лекция 13.}

Пусть $(X, \mathfrak{M}, \mu)$ и $(Y, \mathfrak{N}, \nu)$ — два пространства с мерами, где $\mu$ и $\nu$ — $\sigma$-конечные и полные меры.
Рассмотрим $X \times Y \supset A \times B$, где $A \in \mathfrak{M}$ и $B \in \mathfrak{N}$.
Также рассмотрим
\[\mathcal{P}:=\{A \times B \ | \ A \in \mathfrak{M}, \ B \in \mathfrak{N}, \ \mu(A) < +\infty, \ \nu(B) < +\infty\}\]
— хотим доказать, что это полукольцо.
Будем $A \times B$ называть обобщёнными прямоугольниками.

\begin{lemma}
    $A, A' \subset X$, $B,B' \in Y$, $\{B_\omega\}_{\omega \in I} \subset 2^X$. Тогда
    \begin{enumerate}
        \item \[A \times B \subset A' \times B' \Longleftrightarrow A \subset A' \text{ и } B \subset B',\]
        \item \[(A \times B) \cap (A' \times B') = (A \cap A') \times (B \cap B'),\]
        \item \[(A \times B) \setminus (A' \times B') = (A \setminus A') \times (B \setminus B'),\]
        \item \[A \times \left(\bigcup\limits_{\omega \in I} B_\omega\right) = \bigcup\limits_{\omega \in I} A \times B_\omega,\]
        \item \[A \times \left(\bigcap\limits_{\omega \in I} B_\omega\right) = \bigcap\limits_{\omega \in I} A \times B_\omega.\]
    \end{enumerate}
\end{lemma}
\begin{proof}
    Без доказательства.
\end{proof} 

\begin{theorem}
    Пусть выполнены условия выше. Тогда $\mathcal{P}$ — полукольцо.
    При этом, $\mu \otimes \nu$ — счётно-аддитивная $\sigma$-конечная мера на $\mathcal{P}$ и
    \[\mu \otimes \nu(A \times B) := \mu(A) \nu(B).\]
\end{theorem} 
\begin{proof}
    Поскольку $\mathfrak{M}$ и $\mathfrak{N}$ — $\sigma$-алгебры, следовательно, они полукольца.
    \begin{enumerate}
        \item В начале семестра \hyperref[th:1.1]{было доказано}, что произведение полуколец является полукольцом.
        \item Покажем, что $\mu \otimes \nu$ — счётно-аддитивная мера на $\mathcal{P}$:
        рассмотрим $A \times B \in \mathcal{P}$.
        Пусть 
        \begin{equation}
            \underbrace{A}_{\in \mathfrak{M}} \times \underbrace{B}_{\in \mathfrak{N}} = \bigsqcup\limits_{k=1}^{\infty} \underbrace{A_k}_{\in \mathfrak{M}} \times \underbrace{B_k}_{\in \mathfrak{N}} \label{eq:13.1}
        \end{equation}
        Рассмотрим
        \[\chi_{A \times B} (x,y) = \chi_A(x) \cdot \chi_B(y).\]
        Из \eqref{eq:13.1} следует, что
        \[\chi_{A \times B} (x,y) = \sum\limits_{k=1}^{\infty} \chi_{A_k \times B_k} (x,y) = \sum\limits_{k=1}^{\infty} \chi_{A_k} (x) \cdot \chi_{B_k}(y).\]
        При каждом фиксированном $x \in X$ проинтегрируем полученное равенство по $Y$:
        \[I_x = \int_Y \chi_{A \times B} (x,y) d\nu(y)
        = \chi_A(x) \cdot \int_Y \chi_B(y) d\nu(y)
        = \chi_A(x) \cdot \nu(B).\]
        С другой стороны,
        \begin{multline*}
            I_x = \int_Y \sum\limits_{k=1}^{\infty} \chi_{A_k}(x) \chi_{B_k}(y) d\nu(y)
            = \text{(по теореме Беппо Леви)}
            =\\=
            \sum\limits_{k=1}^{\infty} \int_Y \chi_{A_k}(x) \cdot \int_Y \chi_{B_k}(y) d\nu(y)
            = \sum\limits_{k=1}^{\infty} \chi_{A_k} (x) \nu(B_k)
        \end{multline*} 
        — измеримая неотрицательная функция как ряд из измеримых неотрицательных функций.
        
        Теперь интегрируем по $X$:
        \[\int_X I_x d\mu(x) = \mu(A) \nu(B) = (*)\]
        С другой стороны,
        \begin{multline*}
            \int_X I_x d\mu(x)
            = \int_X \sum\limits_{k=1}^{\infty} \chi_{A_k}(x) \nu(B_k) d\mu(x)
            = \text{(по теореме Беппо Леви)}
            =\\=
            \sum\limits_{k=1}^{\infty} \int_X \chi_{A_k} \nu(B_k)d \mu(x)
            = \sum\limits_{k=1}^{\infty} \mu(A_k) \nu(B_k)
            = \sum\limits_{k=1}^{\infty} \mu \otimes \nu(A_k \times B_k) = (*),
        \end{multline*} 
        следовательно, доказана счётная аддитивность $\mu \otimes \nu$.
        \item Покажем, что $\mu \otimes \nu$ — $\sigma$-конечная мера:
        $\mu$ — $\sigma$-конечная мера, следовательно,
        \[X = \bigcup\limits_{k=1}^{\infty} X_k, \tab \mu(X_k) < +\infty \ \forall k \in \N.\]
        Аналогично $\nu$ — $\sigma$-конечная мера, следовательно,
        \[Y = \bigcup\limits_{j=1}^{\infty} Y_j, \tab \nu(Y_j) < +\infty \ \forall j \in \N,\]
        тогда
        \[X \times Y = \bigcup\limits_{k=1}^{\infty} X_k \times Y = \bigcup\limits_{k=1}^{\infty} \bigcup\limits_{j=1}^{\infty} X_k \times Y_j\]
        — элемент из $\mathcal{P}$ для любой пары $(k,j) \in \N^2$.
    \end{enumerate}
\end{proof} 

Теперь, имея $\sigma$-конечную счётно-аддитивную меру $\mu \otimes \nu$ на полукольце «прямоугольников», построим продолжение $\mu \otimes \nu$ по Каратеодори.
Тогда $\mathfrak{M} \otimes \mathfrak{N}$ — система всех $\mu \otimes \nu$ измеримых по Каратеодори множеств (отметим, что оно шире $\mathcal{P}$).

\begin{remark}
    Если $A$ и $B$ имеют конечную меру, то
    \[A \times B \in \mathfrak{M} \otimes \mathfrak{N} \tab \forall A \in \mathfrak{M}, \ \forall B \in \mathfrak{N},\]
    поскольку $\mathcal{P} \subset \mathfrak{M} \otimes \mathfrak{N}$.
\end{remark}

\begin{remark}
    Пусть теперь $\mu(A) < +\infty$, $\nu(B) = +\infty$, тогда, пользуясь $\sigma$-конечностью меры, разобьём $B$:
    \[B = \bigcup\limits_{k=1}^{\infty} B_k, \tab \nu(B_k) < +\infty \ \forall k \in \N.\]
    \[A \times B = \bigcup\limits_{k=1}^{\infty} (\underbrace{A_k \times B_k}_{\in \mathcal{P}}) \in \mathfrak{M} \otimes \mathfrak{N},\]
    поскольку счётное объединение не выбрасывает за пределы $\sigma$-алгебры.
\end{remark}

\begin{remark}
    Если $\mu(e)=0$, то $\mu \otimes \nu(e \times B) = 0$ для любого $B \in \mathfrak{N}$ — если $\nu(B) < +\infty$, то это очевидно, иначе $B = \bigcup\limits_{k=1}^{\infty} B_k$, $\nu(B_k) < +\infty$, следовательно, в силу счётной-полуаддитивности
    \[\mu \otimes \nu(e \times B) \leqslant \sum\limits_{k=1}^{\infty} \mu \otimes \nu(e \times B_k) = \sum\limits_{k=1}^{\infty} 0 = 0.\]
\end{remark}

Определение произведения двух пространств с мерой обобщается на случай произвольного числа сомножителей.
Например, для трёх пространств
\[(X, \mathfrak{M}, \mu), \tab (Y, \mathfrak{N}, \nu), \tab (Z, \mathfrak{T}, \tau)\]
определяется произведение мер
\[\mu \otimes \nu \otimes \tau = \mu(A) \nu(B) \tau(C),\]
где $A \in \mathfrak{M}$, $B \in \mathfrak{N}$, $C \in \mathfrak{T}$.

\noindent\textbf{Упражнение.} Проверить, что произведение мер ассоциативно.

\subsection{Принцип Кавальери}
Пусть снова
$(X, \mathfrak{M}, \mu)$ — пространство с $\sigma$-конечной полной мерой $\mu$, $(Y, \mathfrak{N}, \nu)$ — пространство с $\sigma$-конечной полной мерой $\nu$.

Для любого $C \subset X \times Y$ введём
\[C_x := \{y \in Y \ | \ (x,y) \in C\},\]
\[C^y := \{x \in X \ | \ (x,y) \in C\}.\]

\begin{lemma}
    $X$,$Y$ — абстрактные множества, $A \subset X \times Y$, $B \subset X \times Y$, $\{B_\omega\} \subset 2^{X \times Y}$.
    \begin{enumerate}
        \item \[(A \setminus B)_x = A_x \setminus B_x \tab \forall x \in X,\]
        \item \[\left(\bigcup\limits_{\omega \in I}B_\omega\right)_x = \bigcup\limits_{\omega \in I}(B_\omega)_x,\]
        \item \[\left(\bigcap\limits_{\omega \in I}B_\omega\right)_x = \bigcap\limits_{\omega \in I} (B_\omega)_x,\]
    \end{enumerate}
    аналогичные утверждения верны для $y$.
\end{lemma} 

\begin{theorem}[Принцип Кавальери]
    Пусть выполнено условие выше. Тогда для любого $C \in \mathfrak{M} \otimes \mathfrak{N}$ выполнено
    \begin{enumerate}
        \item Для $\mu$-почти всех $x \in X$ верно $C_x \in \mathfrak{N}$;
        \item $x \to \nu(C_x)$ — измеримая в обобщённом смысле функция (то есть эта функция измерима на множестве $X_0 \subset X$: $\mu(X \setminus X_0) = 0$);
        \item $$\mu \otimes \nu(C) = \int_{X_0} \nu(C_x) d\mu(x).$$
    \end{enumerate}
    аналогичные утверждения верны по переменной $y$.
\end{theorem} 
\begin{proof}
    Доказательство проведём в 4 шага. На первых трёх шагах предположим, что меры конечны. На 4 шаге избавимся от этого предположения.
    \begin{enumerate}
        \item Рассмотрим систему $\mathfrak{R}$ всех таких $C \in \mathfrak{M} \otimes \mathfrak{N}$, для которых $C_x$ измерима при всех $x \in X$ и $x \to \nu(C_x)$ — измеримая функция (в обычном смысле).
        
        Заметим, что $\mathfrak{R} \supset \mathcal{P}$.
        Кроме того, система $\mathfrak{R}$ является монотонным классом.

        Действительно, пусть $C_1 \subset \dots \subset C_n \subset \dots$, $C = \bigcup\limits_{k=1}^{\infty}C_k$ — возрастающая последовательность множеств из $\mathfrak{R}$, следовательно, для любого $x \in X$
        \[C_x = \bigcup\limits_{k=1}^{\infty}(\underbrace{C_k}_{\in \mathfrak{N}})_x \implies C_x \in \mathfrak{N}.\]
        Вместе с тем,
        \[\nu(C_x) = \text{(в силу непрерывности меры снизу)} = \lim\limits_{k\to \infty}\nu((C_k)_x).\]
        Определим
        \[f_k(x):=\nu((C_k)_x), \tab f_{k+1}(x) \geqslant f_k(x)\]
        — измеримая функция от $x$ для любого $k \in \N$.
        Также определим предельную функцию:
        \[f(x):=\nu(C_x) = \lim\limits_{k\to \infty}f_k(x)\]
        — измеримая функция на $X$ как предел.
        Следовательно, $C \in \mathfrak{R}$.

        Заметим, что для любого $A \in \mathfrak{R}$ дополнение $A^c$ также принадлежит $\mathfrak{R}$.
        \[(A^c)_x = Y \setminus A_x \in \mathfrak{N}, \tab \forall x \in X.\]
        \[\nu((A^c)_x) = \nu(Y) - \nu(A_x), \tab \forall x \in X\]
        — именно здесь используется конечность меры.
        Следовательно, $x \to \nu((A^c)_x)$ — измеримая функция от $x$.
        Тогда $\forall B_1 \supset \dots \supset B_n \supset \dots$, $B_n \in \mathfrak{R}$ верно $\bigcap\limits_{n=1}^{\infty}B_n \in \mathfrak{R}$.

        Теперь покажем, что $\mathfrak{R}$ — алгебра.
        Если $A \in \mathfrak{N}$, $B \in \mathfrak{N}$, $A \cap B = \emptyset$, то для любого $x \in X$
        \[(A \sqcup B)_x = A_x \sqcup B_x \in \mathfrak{N}\]
        и
        \[\nu((A \sqcup B)_x) = \nu(A_x) + \nu(B_x),\]
        тогда
        $x \to \nu((A \sqcup B)_x)$ — измеримая функция на $X$.
        В силу \hyperref[th:1.2]{теоремы о дизъюнктном представлении в полукольце} если $C = \bigcup\limits_{k=1}^{N} C_k$, где $C_1, \dots, C_N \in \mathcal{P}$, то существует дизъюнктный набор $\{B_j\}_{j=1}^L \subset \mathcal{P}$: $C = \bigsqcup\limits_{j=1}^{L}B_j$, следовательно, $C \in \mathfrak{R}$ ($\mathfrak{R}$ замкнуто относительно дизъюнктного объединения).

        Таким образом, мы поняли, что $\mathfrak{R}$ — монотонный класс, $\mathfrak{R} \supset \mathcal{P}$, $\mathfrak{R}$ замкнуто относительно конечного объединения и относительно дополнения, следовательно, $\mathfrak{R}$ — алгебра и монотонный класс, а потому $\mathfrak{R}$ содержит минимальную алгебру полукольца $\mathcal{P}$, то есть $\mathfrak{R} \supset \mathfrak{M}(\mathcal{P})$.
        Условия 1-2 теоремы выполнены для любых $C \in \mathfrak{M}(\mathcal{P})$.

        Теперь покажем, что
        \[\mu \otimes \nu(C) = \int_X \nu(C_x) d\mu(x) \tab \forall C \in \mathfrak{M}(\mathcal{P}).\]
        Рассмотрим меру $\Phi$ на $\mathfrak{M}(\mathcal{P})$:
        \[\underbrace{E}_{\in \mathfrak{M}(\mathcal{P})} \xrightarrow{\Phi} \int_X \nu(E_x)d\mu(x).\]
        Покажем, что $\Phi$ — это мера.
        \begin{multline*}
            \Phi \left(\bigsqcup\limits_{n=1}^{\infty}E_n\right)
            =
            \int_X \nu \left(\left(\bigsqcup\limits_{n=1}^{\infty}E_n\right)_x\right) d\mu(x)
            =
            \int_X \sum\limits_{n=1}^{\infty} \nu((E_n)_x) d\mu(x)
            =\\=
            \text{(по теореме Беппо Леви)}
            =
            \sum\limits_{n=1}^{\infty} \int_X \nu((E_n)_x) d\mu(x)
            =
            \sum\limits_{n=1}^{\infty}\Phi(E_n) \implies
        \end{multline*} 
        $\Phi$ — мера на $\mathfrak{M} \otimes \mathfrak{N}$.
        Но $\Phi|_\mathcal{P} = \mu \otimes \nu|_\mathcal{P}$, следовательно, $\Phi = \mu \otimes \nu$ на $\mathfrak{M}(\mathcal{P})$ по теореме о единственности продолжения меры.
        \item Рассмотрим множество 
        $$e \in \mathfrak{M} \otimes \mathfrak{N}: \ \mu \otimes \nu(e) = 0.$$
        \[\exists \widetilde{e} \in \mathfrak{M}(\mathcal{P}): \mu \otimes \nu(\widetilde{e}) = 0.\]
        Для $\widetilde{e}$ выполнены условия 1,2,3, то есть
        \begin{itemize}
            \item $(\widetilde{e})_x \in \mathfrak{N}$ для любого $x \in X$;
            \item $x \to \nu(\widetilde{e}_x)$ измерима;
            \item $\mu \otimes \nu(\widetilde{e}) = \int_X \nu(\widetilde{e}_x)d\mu(x) = 0,$ 
        \end{itemize}
        следовательно, $\nu(\widetilde{e}_x) = 0$ для $\mu$-почти всех $x \in X$. Обозначим $X_0$ — множество, где выполнено данное условие.
        Поскольку $\nu$ — полная мера, $e_x \subset \widetilde{e}_x$, следовательно, $e_x$ измеримо для $\mu$-почти всех $x \in X$ и $\nu(e_x) = 0$ $\mu$-почти всюду.
        Тогда $x \to \nu(e_x)$ измерима в широком смысле на $X$.
        Получается,
        \[\mu \otimes \nu(e) = \int_X \nu(e_x) d\mu(x).\]
        \item
        \[\forall C \in \mathfrak{M} \otimes \mathfrak{N} \ \exists \widetilde{C} \in \mathfrak{M}(\mathcal{P}): \ C \subset \widetilde{C} \text{ и } \mu \otimes \nu(C) = \mu \otimes \nu(\widetilde{C}) \implies \widetilde{C} = C \cup e.\]
        В силу предыдущих шагов $C_x \in \mathfrak{N}$ $\mu$-почти всюду (т.е. всюду на $X_0$) и $x \to \nu(C_x)$ измеримо в широком смысле на $X$ (в обычном смысле на $X$).
        \begin{multline*}
            \mu \otimes \nu(C)
            =
            \mu \otimes \nu(\widetilde{C})
            =
            \int_X \nu(\widetilde{C}_x) d\mu(x)
            =
            \int_X \nu(C_x)d\mu(x)
            =
            \int_{X_0}\nu(C_x)d\mu(x)
        \end{multline*} 
        \item Отбрасываем требование конечности меры, используя $\sigma$-конечность меры:
        \[X = \bigsqcup\limits_{k=1}^{\infty}X_k, \tab \mu(X_k) < +\infty \ \forall k \in \N, \ X_k \in \mathfrak{M},\]
        \[Y = \bigsqcup\limits_{n=1}^{\infty}Y_n, \tab \nu(Y_n) < +\infty \ \forall n \in \N, \ Y_n \in \mathfrak{N}.\]
        Пусть $C \in \mathfrak{M} \otimes \mathfrak{R}$, тогда
        \[C = \bigsqcup\limits_{n,k}C_{k,n},\]
        где $C_{k,n}:=C \cap (X_k \times Y_n)$.
        Применим предыдущие шаги к $X = X_k$ и $Y = Y_n$ для всех $k,n$, $\mathfrak{M}|_{X_k}$ и $\mathfrak{N}|_{Y_k}$ и получим, что для всех $k,n \in \N$ $(C_{k,n})_x \in \mathfrak{N}$ $\mu$-почти всюду, а $x \to \nu((C_{k,n})_x)$ измерима в широком смысле на $X_k$ (следовательно, и на $X$ при доопределении тождественным нулём за пределами $X_k$).
        \[\mu \otimes \nu(C_{k,n}) = \int_{X_k} \nu((C_{k,n})_x)d\mu(x).\]
        Теперь соберём всё вместе:
        \[C_x = \bigsqcup\limits_{k} \bigsqcup\limits_{n} (C_{k,n})_x\]
        — измеримо $\mu$-почти всюду на $X$.
        \[\nu(C_x) = \nu \left(\bigsqcup\limits_{k} \bigsqcup\limits_{n}(C_{k,n})_x\right) = \sum\limits_{k} \sum\limits_{n} \nu((C_{k,n})_x)\]
        и $x \to \nu(C_x)$ измерима в широком смысле как ряд из неотрицательных измеримых в широком смысле функций $x \to \nu((C_{k,n})_x)$.

        Покажем, что справедливо условие 3.
        \begin{multline*}
            \int_X \nu(C_x)d\mu(x)
            =
            \int_X \nu \left(\left(\bigsqcup\limits_{k,n=1}^{\infty}(C_{k,n})\right)_x\right)d\mu(x)
            =\\=
            \int_X \sum\limits_{k=1}^{\infty}\nu\left(\bigsqcup\limits_{n=1}^{\infty}(C_{k,n})_x\right)d\mu(x)
            =
            \text{(по теореме Беппо Леви)}
            =\\=
            \sum\limits_{k=1}^{\infty}\int_X \nu\left(\bigsqcup\limits_{n=1}^{\infty}(C_{k,n})_x\right)d\mu(x)
            =
            \sum\limits_{k=1}^{\infty}\int_X \sum\limits_{n=1}^{\infty} \nu((C_{k,n})_x)d\mu(x)
            =\\=
            \sum\limits_{k=1}^{\infty} \sum\limits_{n=1}^{\infty} \int_{X_k} \nu((C_{k,n})_x)d\mu(x)
            =
            \sum\limits_{k,n=1}^{\infty}\mu \otimes \nu(C_{k,n})
            =
            \mu(C). 
        \end{multline*} 
    \end{enumerate}
\end{proof} 


\paragraph{Лекция 14.}

\begin{theorem}[Тонелли]
    Пусть $(X, \mathfrak{M}, \mu)$ и $(Y, \mathfrak{N}, \nu)$ — пространства с $\sigma$-конечными и полными мерами.
    Пусть $f: X \times Y \to [0, +\infty]$ — измеримая относительно $\mathfrak{M} \otimes \mathfrak{N}$ функция.
    Тогда:
    \begin{enumerate}
        \item Для $\mu$-почти всех $x \in X$ $f_x(\cdot)$ измерима на $Y$:
        \[f_x(y):=f(x,y), \tab \forall y \in Y;\]
        \item $x \to \int_X f_x(y) d\nu(y)$ измерима в широком смысле как функция на $X$;
        \item \[\int_{X \times Y} f(x,y)d(\mu \otimes \nu)(x,y) = \int_X \left(\int_Y f(x,y)d\nu(y)\right)d\mu(x).\]
    \end{enumerate}
    Аналогичные утверждения справедливы при замене ролей $X$ и $Y$.
\end{theorem} 
\begin{proof}
\begin{enumerate}
    \item Пусть $f = \chi_C$, где $C \in \mathfrak{M} \otimes \mathfrak{N}$.
    \[\chi_C(x,y) = \chi_{C_x}(y) = f_x(y) = 
    \begin{cases}
        1, \ (x,y) \in C,\\
        0, \ \text{иначе}.
    \end{cases}
    \]
    В силу принципа Кавальери с прошлой лекции все пункты теоремы верны.
    \item Пусть $f$ — простая функция.
    \[f = \sum\limits_{k=1}^{N} a_k \chi_{C_k}, \ \{C_i\}_{i=1}^N \subset \mathfrak{M} \otimes \mathfrak{N}, \ C_i \cap C_j = \emptyset \text{ при } i \neq j, \ a_k \in \R.\]
    \[f_x = \sum\limits_{k=1}^{N}a_k (\chi_{C_k})_x.\]
    Поскольку $C_k$ — конечное число, то из предыдущего шага легко видеть, что условия теоремы выполнены.
    \item В силу теоремы Рисса имеется поточечая сходимость:
    \[f = \lim\limits_{n\to \infty}f_n, \ \{f_n\} \text{ — монотонная последовательность простых функций.}\]
    Для любого $n \in \N$ существует $E_n \subset X$: $\mu(E_n)=0$, следовательно, для любого $x \in X \setminus E_n$ $(f_n)_x$ измерима на $Y$ и $x \to \int_X (f_n)_x d\nu$ измерима на $X \setminus E_n$.
    
    Определим
    \[\underline{X} := X \setminus \bigcup\limits_{n=1}^{\infty}E_n,\]
    тогда для любого $x \in \underline{X}$ и для любого $n \in \N$ $(f_n)_x$ — измеримая на $Y$ функция, а $x \xrightarrow{\Phi_n} \int_X (f_n)_x d\nu$ — измеримая на $\underline{X}$ функция.

    Очевидно, что для любого $x \in X$ $f(x,y)=\lim\limits_{n\to \infty}f_n(x,y)$ и 
    \[f_x(y)=\lim\limits_{n\to \infty}(f_n)_x (y), \ \forall y \in Y.\]
    В силу предыдущего рассуждения (поскольку поточечный предел не портит измеримость) $f_x$ измерима на $Y$ для любого $x \in \underline{X}$.
    В силу монотонности последовательности
    \[f_{n+1}(x,y) \geqslant f_n(x,y) \ \forall x,y \in X \times Y, \ \forall n \in \N \implies\]
    \[\forall x \in X: \ (f_{n+1})_x(y) \geqslant (f_n)_x(y) \ (*) \ \forall n \in \N \text{ и } \forall y \in Y \implies\]
    в силу теоремы Беппо Леви:
    \[\lim\limits_{n\to \infty} \underbrace{\int_Y (f_n)_x(y) d\nu(y)}_{=:\phi_n(x)} = \underbrace{\int_Y (f)_x (y) d\nu(y)}_{=: \phi(x)} \ \forall x \in \underline{X}
    \implies
    \phi_n \xrightarrow{\underline{X}} \phi, \ n \to \infty,\]
    но $\phi_n$ измерима на $\underline{X}$, следовательно, $\phi$ измерима на $\underline{X}$, то есть $\phi$ измерима в широком смысле на $X$.

    Покажем, что выполнимо равенство 3:
    \begin{multline*}
    \int_{X \times Y} f(x,y) d(\mu \otimes \nu)(x,y)
    =\\=
    \lim\limits_{n\to \infty} \int_{X \times Y} f_n(x,y) d(\mu \times \nu)(x,y)
    =\\=
    \lim\limits_{n\to \infty} \int_X (\phi_n)_x(y) d\nu(y) =(**)
    \end{multline*} 
    из $(*)$ следует, что
    \[\phi_{n+1}(x) \geqslant \phi_n(x) \ \forall x \in \underline{X}.\]
    \begin{multline*}
        (**) = \lim\limits_{n\to \infty} \int_{\underline{X}} (\phi_n)_x (y) d\nu(y)
        =
        \int_{\underline{X}} \lim\limits_{n\to \infty} (\phi_n)_x (y) d\nu(y)
        =
        \int_X \phi(y) d\nu(y).
    \end{multline*} 
\end{enumerate}
\end{proof} 

\begin{theorem}[Фубини]
    Пусть $(X, \mathfrak{M}, \mu)$ и $(Y, \mathfrak{N}, \nu)$ — пространства с полными $\sigma$-конечными мерами.
    Пусть $f \in \widetilde{L}_1(X \times Y, \mathfrak{M} \times \mathfrak{N}, \mu \otimes \nu)$. Тогда
    \begin{enumerate}
        \item Для $\mu$-почти всех $x \in X$ $f_x \in \widetilde{L}_1(Y, \mathfrak{N}, \nu)$;
        \item \[\phi(x) = \int_Y f_x(y) d\nu(y), \ \phi \in \widetilde{L}_1(X, \mathfrak{M}, \mu);\]
        \item \[\int_{X \times Y} f(x,y) d(\mu \otimes \nu)(x,y) = \int_X \left(\int_Y f_x(y) d\nu(y)\right)d\mu(x).\]
    \end{enumerate}
    Аналогичное утверждение справедливо при замене $(X, \mathfrak{M}, \mu)$ и $(Y, \mathfrak{N}, \nu)$.
\end{theorem} 
\begin{proof}
    Напомним, что $f = f_+ - f_-$.
    К $f_+$ и $f_-$ применим теорему Тонелли:
    \[+\infty > \int_{X \times Y} f_{\pm}(x,y)d(\mu \otimes \nu)(x,y) = \int_X \left(\int_Y (f_{\pm})_x (y) d\nu(y)\right)d\mu(x).\]
    В силу свойств интеграла Лебега для $\mu$-почти всех $x \in X$
    \[\int_Y (f_{\pm})_x(y)d\nu(y) < +\infty,\]
    следовательно, для $\mu$-почти всех $x \in X$ $f_x = (f_+)_x - (f_-)_x \in \widetilde{L}_1(X, \mathfrak{M}, \mu)$.
    \begin{multline*}
        \int_{X \times Y} f(x,y) d(\mu \otimes \nu)(x,y)
        =\\=
        \int_{X \times Y} f_+(x,y)d(\mu \otimes \nu)(x,y) - \int_{X \times Y} f_-(x,y) d(\mu \otimes \nu)(x,y)
        =\\=
        \int_X \left(\int_Y (f_+)_x(y)d\nu(y)\right)d\mu(x) - \int_X \left(\int_Y (f_-)_x(y)d\nu(y)\right)d\mu(x) = (*)
    \end{multline*} 
    \[\phi_+(x) = \int_Y (f_+)_x d\nu, \ \phi_-(x) = \int_X (f_-)_x d\nu, \ \phi_+, \ \phi_- \in \widetilde{L}_1(X, \mathfrak{M}, \mu).\]
    \[(*) = \int_X (\phi_+(x) - \phi_-(x)) d\mu(x).\]
\end{proof} 

\begin{example}
    \[Q = [-1,1]^2,\]
    \[f(x,y) = 
    \begin{cases}
        \frac{x^2-y^2}{(x^2+y^2)^2}, \ x^2+y^2>0 \\
        0, \ \text{иначе}.
    \end{cases}
    \]
    \[\int_{-1}^1 \left(\int_{-1}^1 f(x,y)dx\right)dy = -\pi \neq \pi = \int_{-1}^1 \left(\int_{-1}^1 f(x,y)dy\right)dx,\]
    \[f \notin \widetilde{L}_1(Q).\]
\end{example}

\begin{example}
    \[\int_{-1}^1 \int_{-1}^1 \frac{xy}{(x^2+y^2)^2} dxdy = 0 = \int_{-1}^1 \int_{-1}^1 \frac{xy}{(x^2+y^2)^2}dydx,\]
    но эта функция не интегрируема на $Q$.
\end{example}
